\section{Chain, product, and quotient rules}

Most functions are not single elementary building blocks. They are combinations of simpler functions. To differentiate them, we need rules that describe how derivatives behave under products, quotients, and composition.

The first task is always to read the structure of the expression.

\subsection*{The product rule}

The derivative of a product is not the product of the derivatives. Both factors may change, and the rule must account for both changes.

\begin{theorembox}
If \(f\) and \(g\) are differentiable, then
\[
(fg)'=f'g+fg'.
\]
Equivalently,
\[
\frac{d}{dx}\bigl(f(x)g(x)\bigr)=f'(x)g(x)+f(x)g'(x).
\]
\end{theorembox}

\begin{examplebox}
Differentiate
\[
y=x^2e^x.
\]
Let
\[
f(x)=x^2,
\qquad
 g(x)=e^x.
\]
Then
\[
f'(x)=2x,
\qquad
 g'(x)=e^x.
\]
By the product rule,
\[
y'=2x e^x+x^2 e^x.
\]
We may factor:
\[
y'=e^x(2x+x^2).
\]
\end{examplebox}

\subsection*{The quotient rule}

For a quotient, the denominator also changes. The quotient rule is:

\begin{theorembox}
If \(f\) and \(g\) are differentiable and \(g(x)\ne0\), then
\[
\left(\frac{f}{g}\right)'=\frac{f'g-fg'}{g^2}.
\]
\end{theorembox}

\begin{examplebox}
Differentiate
\[
y=\frac{x^2+1}{x-1}.
\]
Let
\[
f(x)=x^2+1,
\qquad
 g(x)=x-1.
\]
Then
\[
f'(x)=2x,
\qquad
 g'(x)=1.
\]
By the quotient rule,
\[
y'=\frac{2x(x-1)-(x^2+1)\cdot1}{(x-1)^2}.
\]
Simplify the numerator:
\[
2x(x-1)-(x^2+1)=2x^2-2x-x^2-1=x^2-2x-1.
\]
Thus
\[
y'=\frac{x^2-2x-1}{(x-1)^2}.
\]
\end{examplebox}

The quotient rule should be used with attention to the domain. The original function is defined only where the denominator is not zero.

\subsection*{The chain rule}

The chain rule is used for compositions of functions. If an expression has an inside function and an outside function, the derivative must account for both.

\begin{theorembox}
If
\[
y=f(g(x)),
\]
then
\[
y'=f'(g(x))\,g'(x).
\]
\end{theorembox}

A useful way to say this is: differentiate the outside function, keep the inside function, and multiply by the derivative of the inside function.

\begin{examplebox}
Differentiate
\[
y=(3x+1)^7.
\]
The outside function is the seventh power, and the inside function is
\[
3x+1.
\]
By the chain rule,
\[
y'=7(3x+1)^6\cdot3.
\]
Therefore
\[
y'=21(3x+1)^6.
\]
\end{examplebox}

\begin{examplebox}
Differentiate
\[
y=\sin(x^2).
\]
The outside function is \(\sin u\), and the inside function is \(u=x^2\). Thus
\[
y'=\cos(x^2)\cdot 2x.
\]
So
\[
y'=2x\cos(x^2).
\]
\end{examplebox}

\subsection*{Reading nested expressions}

The chain rule is often the most important rule because many expressions are nested. Before differentiating, identify the layers.

\begin{examplebox}
Differentiate
\[
y=e^{\cos^2 x}.
\]
This expression has several layers:
\[
x \longmapsto \cos x \longmapsto (\cos x)^2 \longmapsto e^{\cos^2 x}.
\]
Let
\[
u=\cos^2 x.
\]
Then
\[
y=e^u,
\qquad
 y'=e^u u'.
\]
Now
\[
u=(\cos x)^2,
\]
so by the chain rule again,
\[
u'=2\cos x\cdot(-\sin x).
\]
Therefore
\[
y'=e^{\cos^2 x}\cdot 2\cos x(-\sin x).
\]
Thus
\[
y'=-2e^{\cos^2 x}\sin x\cos x.
\]
\end{examplebox}

\begin{summarybox}
Before differentiating a combined expression, ask:
\begin{itemize}
\item Is the expression a sum, product, quotient, or composition?
\item Are there nested functions?
\item What is the inside function?
\item What is the outside function?
\item Do I need more than one rule?
\item Can the result be simplified without hiding its structure?
\end{itemize}
\end{summarybox}

Most differentiation errors come from misreading the structure of the expression. Read first, then apply the rule.